% Part: sets-functions-relations
% Chapter: arithmetization
% Section: checking-details
%
\documentclass[../../../include/open-logic-section]{subfiles}

\begin{document}

\olfileid[ar]{sfr}{arith}{check}
\olsection{الحلقات والحقول المرتبة}

تكرر في الفصل قولنا إن التعريفات تسلك «كما ينبغي». فنفصل في هذا الملحق التقني المراد من ذلك، ونبين على وجه الإجمال طريق البرهان على سلامة هذا السلوك.

قد عرفنا الجمع والضرب على $\Int$ في \olref[int]{sec}. ومطلبنا أن تمنح هاتان العمليتان $\Int$ البنية التي قصدناها لها، وهي على التعيين بنية الحلقة التبديلية.

\begin{defn}
\emph{الحلقة التبديلية} مجموعة $S$ عين فيها عنصران $0$ و$1$ وعمليتان $+$ و$\times$، واستوفت الصيغ الثماني:
\begin{align*}
		\emph{التجميعية}&&a + (b+ c) & = (a + b) + c \\
		&& (a \times b) \times c & = a \times (b\times c)\\
		\emph{التبديلية}&&a + b &= b+ a  \\
		&&  a \times b&= b\times a\\ 
		\emph{العنصران المحايدان}&&a + 0 &= a \\
		&& a \times 1 &= a\\
		\emph{المعكوس الجمعي}&&(\exists b\in S)0&=a + b\\
		\emph{التوزيعية}&&a \times (b+ c ) &= (a \times b) + (a \times c)
	\end{align*}
	ونفهم هذه الصيغ جميعًا مكممة كليًّا، مع تقييد المتغيرات بـ$S$. وليس من شرط العنصرين $0$ و~$1$ أن يكونا العددين الطبيعيين المعروفين بهذين الاسمين.
\end{defn}

فإثبات أن الأعداد الصحيحة حلقة تبديلية يرجع إلى تحقيق هذه الشروط الثمانية. وليس واحد منها {صعبًا}، وإن كان في استيفائها شيء من المشقة. ونمثل لذلك ببرهان \emph{التجميعية} للجمع.

\begin{proof}
خذ $i, j, k \in \Int$. فتوجد $a_1, b_1, a_2, b_2, a_3, b_3 \in
\Nat$ تحقق $i = \equivrep{a_1, b_1}{}$ و$j =
\equivrep{a_2,b_2}{}$ و$k = \equivrep{a_3, b_3}{}$. ونختصر في هذا القسم كله بكتابة «$\equivrep{x, y}{}$» مكان «$\equivrep{x, y}{\Intequiv}$» تيسيرًا للقراءة. فنحسب
\begin{align*}
	i + (j + k) &= \equivrep{a_1, b_1}{}+(\equivrep{a_2, b_2}{} + \equivrep{a_3, b_3}{}) \\
	&= \equivrep{a_1,  b_1}{} + \equivrep{a_2+a_3, b_2+b_3}{}\\
	&= \equivrep{a_1 + (a_2 + a_3), b_1 + (b_2 + b_3)}{}\\
	&= \equivrep{(a_1 + a_2) + a_3, (b_1 + b_2) + b_3}{}\\
	&= \equivrep{a_1 + a_2, b_1 + b_2}{} + \equivrep{a_3, b_3}{}\\
	&= (\equivrep{a_1, b_1}{} + \equivrep{a_2, b_2}{}) + \equivrep{a_3, b_3}{}\\
	&= (i+j) + k
\end{align*}
والاعتماد في هذه التحويلات على قوانين الجمع في $\Nat$.
\end{proof}

ونبين على النهج نفسه \emph{وجود المعكوس الجمعي}.

\begin{proof}
خذ $i \in \Int$، واكتب $i = \equivrep{a,b}{}$ لبعض $a,b \in
\Nat$. ثم ضع $j = \equivrep{b,a}{} \in \Int$. فمن قوانين الطبيعية $(a+b) + 0 = 0 + (a+b)$؛ فيلزم بالتعريف $\tuple{a+b, b+a} \sim_\Int \tuple{0,0}$، ومنه $\equivrep{a+b, b+a}{} = \equivrep{0, 0}{} = 0_\Int$. ولذلك $i + j =
\equivrep{a,b}{}+\equivrep{b,a}{}=\equivrep{a+b, b+a}{}= \equivrep{0,
0}{} = 0_\Int$، وهو المطلوب.
\end{proof}

وأما \emph{التوزيعية} فبرهانها بالحساب الآتي.

\begin{proof}
نأخذ، كما تقدم، $i = \equivrep{a_1, b_1}{}$ و$j =
\equivrep{a_2,b_2}{}$ و$k = \equivrep{a_3, b_3}{}$؛ فنجد:
\begin{align*}
	i \times (j + k) 
	&= \equivrep{a_1, b_1}{} \times (\equivrep{a_2,b_2}{} + \equivrep{a_3, b_3}{})\\
	&= \equivrep{a_1, b_1}{} \times \equivrep{a_2 + a_3,b_2+b_3}{}\\
	&= \equivrep{a_1  (a_2 + a_3) + b_1  (b_2+b_3), a_1  (b_2 + b_3) + b_1 (a_2 + a_3)}{}\\
	&= \equivrep{a_1 a_2 + a_1a_3 + b_1 b_2+b_1b_3, a_1 b_2 + a_1b_3 + a_2b_1 + a_3b_1}{}\\		
%		&= \equivrep{a_1 a_2 + b_1b_2 + a_1a_3 + b_1 b_2+b_1b_3, a_1 b_2 + a_1b_3 + a_2b_1 + a_3b_1}{}\\		
	&= \equivrep{a_1a_2 + b_1b_2, a_1b_2 + a_2b_1}{} + \equivrep{a_1a_3 + b_1b_3, a_1b_3 + a_3b_1}{}\\
	&= (\equivrep{a_1, b_1}{} \times \equivrep{a_2,b_2}{}) + (\equivrep{a_1, b_1}{} \times  \equivrep{a_3, b_3}{})\\
	&= (i \times j) + (i \times k)
\end{align*}
\end{proof}

وأما الشروط الخمسة الباقية فإثباتها تمرين. وباستيفائها يثبت أن $\Int$ حلقة تبديلية؛ أي إن الجمع والضرب اللذين عرفناهما يؤديان ما طلبناه منهما.

\begin{prob}
حقق أن $\Int$ حلقة تبديلية.
\end{prob}

ولا ينتهي العمل بذلك؛ فقد عرفنا على $\Int$، إلى جانب الجمع والضرب، علاقة الترتيب $\leq$. فينبغي تحقيق سلامتها أيضًا، أي إثبات أن $\Int$ حلقة \emph{مرتبة}.\footnote{مر في \olref[sfr][rel][ord]{def:linearorder} أن الترتيب الكلي علاقة انعكاسية متعدية مضادة للتناظر متصلة. ويسمى الاتصال أحيانًا في هذا السياق \emph{التثليث}؛ إذ لكل $a$ و$b$ يتحقق $a \leq b \lor a
= b \lor a \geq b$.}

\begin{defn}
\emph{الحلقة المرتبة} حلقة تبديلية عين فيها ترتيب كلي $\leq$ يحقق:
\begin{align*}
	a \leq b &\lif a + c \leq b + c\\
	(a \leq b \land 0 \leq c) &\lif a \times c \leq b \times c
\end{align*}
\end{defn}

\begin{prob}
حقق أن $\Int$ حلقة مرتبة.
\end{prob}

فإثبات أن $\Int$ على هذا البناء حلقة مرتبة يجري على حساب مألوف، وإن طال بعض الطول؛ ونكل استيفاءه إلى القارئ.

وبهذا يفرغ النظر في الصحيحة، ونطلب للنسبية أحكامًا قريبة منها. والمقصود خاصة أنها \emph{حقل} مرتب، وفق ما عرفناه من $+$ و$\times$ و$\leq$:
\begin{defn}\ollabel{orderedfield}
% تصحيح تعريفي (0047-I1): استُبعدت الحلقة الصفرية بإضافة الشرط القياسي 0 != 1.
\emph{الحقل المرتب} حلقة مرتبة يكون فيها $0 \neq 1$، ويضاف إلى شروطها:
\begin{align*}
	\emph{المعكوس الضربي}& & (\forall a \in S \setminus \{0\})(\exists b \in S) a\times b& = 1
\end{align*}
\end{defn}

\textbf{تنبيه تصحيحي على الأصل.} أُضيف شرط $0 \neq 1$ إلى تعريف
الحقل؛ إذ كانت الشروط المطبوعة تجيز الحلقة ذات العنصر الواحد، وفيها
$0=1$، مع أنها ليست حقلًا بالاصطلاح المقصود هنا.

فبعد أن تثبت كون $\Int$ حلقة مرتبة، يسهل عليك إثبات أن $\Rat$ حقل مرتب، وإن احتاج الإثبات إلى شيء من العمل.

\begin{prob}
حقق أن $\Rat$ حقل مرتب.
\end{prob}

فإذا فرغنا من الصحيحة والنسبية بقيت الحقيقية. ومطلبنا أن $\Real$ حقل مرتب \emph{كامل}؛ أي حقل مرتب ذو خاصية الاكتمال. وقد ثبت في \olref[cuts]{realcompleteness} أن $\Real$ كاملة، وبقي استيفاء الحساب الذي يثبت أن $\Real$ حقل مرتب، وفيه طول.

وقبل الشروع في \emph{ذلك} التمرين، نحقق أمورًا أقرب مأخذًا. فمنها أن $\alpha + \beta$ كما عرفناه \emph{قطع} حقًّا، متى كان $\alpha$ و$\beta$ قطعين. وبيانه الآتي.

\begin{proof}
% استكمال برهاني (0047-I2): أُثبتت حقيقية مجموعة المجاميع بدل افتراضها من كون الحدين قطعين.
لما كان $\alpha$ و$\beta$ قطعين، كانت $\alpha + \beta = \Setabs{p
+ q}{p \in \alpha \land q \in \beta}$ مجموعة غير خالية. واختر
$r\in\Rat\setminus\alpha$ و$s\in\Rat\setminus\beta$؛ فمن ابتدائية
القطعين يكون $p<r$ و$q<s$ لكل $p\in\alpha$ و$q\in\beta$، فيغدو $r+s$
حدًّا أعلى لا يقع في مجموعة المجاميع؛ فهي جزء حقيقي من $\Rat$. وإذا كان
$x < p + q$ لبعض $p \in \alpha$ و$q \in \beta$، حصل $x - p < q$، ومن ابتدائية القطع $x - p \in \beta$؛ فيكون $x = p + (x - p) \in
\alpha + \beta$. فالمجموعة $\alpha + \beta$ مقطع ابتدائي من $\Rat$. وأخيرًا، خذ $p + q \in \alpha + \beta$؛ فبما أن $\alpha$ و$\beta$ قطعان، يوجد $p_1 \in \alpha$ و$q_1 \in \beta$ بحيث $p < p_1$ و$q < q_1$. ومنه $p + q < p_1 + q_1 \in \alpha +
\beta$، فلا عنصر أعظم في $\alpha + \beta$.
\end{proof}

وعلى هذا النحو يمكن تحقيق أن $\alpha - \beta$ و$\alpha
\times \beta$ و$\alpha \div \beta$ قطوع أيضًا؛ مع استثناء كون $\beta$ قطع الصفر في القسمة. ونترك استيفاء هذه الحجج للقارئ.

\begin{prob}
حقق أن $\Real$ حقل مرتب.
\end{prob}

وبقي أمر يسير أرجئ تحقيقه. فقد وضعنا في \olref[cuts]{sec} أن $\sqrt{2} =
\Setabs{p \in \Rat}{p < 0 \text{ أو }p^2 < 2}$، ولا بد من إثبات أن هذه المجموعة \emph{قطع}.

\begin{proof}
% استكمال برهاني (0047-I3): عولجت العناصر غير الموجبة قبل حجة التقريب المخصصة للعناصر الموجبة.
كون هذه المجموعة مقطعًا ابتدائيًّا حقيقيًّا غير خال من النسبية ظاهر. فيبقى انتفاء العنصر الأعظم؛ فإذا كان~$p\leq0$ كان~$p<1$، وكان~$1$ في المجموعة لأن~$1^2<2$. وتبقى حال العدد النسبي الموجب~$p$ الذي يحقق $p^2 < 2$؛ فنضع $q = \frac{2p+2}{p+2}$، ثم نثبت $p < q$ و$q^2 < 2$. أما $p < q$ فمن الحساب
\begin{align*}
	p^2 &< 2\\
	p^2 + 2p &< 2 + 2p\\
	p(p + 2) &< 2 + 2p\\
	p &< \tfrac{2+2p}{p+2} = q
\end{align*}
وأما $q^2 < 2$ فينتج كما يأتي:
\begin{align*}
	p^2 &< 2\\
	2p^2 + 4p + 2 &< p^2 + 4p+ 4\\
	4p^2 + 8p + 4 &< 2(p^2 + 4p + 4)\\
	(2p+2)^2 & <2(p+2)^2\\
	\tfrac{(2p+2)^2}{(p+2)^2} &< 2\\
	q^2 &< 2
\end{align*}
\end{proof}
\end{document}
